% GENERATED FILE. Edit canonical lessons, then run npm run generate:books.
% canonical-source-hash: fnv1a64:dca0e9bb88b13d7f
% canonical-lessons: JA-C216-byouki, JA-C216-guai, JA-C216-tsukareru, JA-C216-naoru, JA-C216-furueru

\chapter{Feeling Unwell}
\label{ch:ja-feeling-unwell}
\hlchaptermodality{216}

\begin{chapteropening}
\textbf{By the end of this chapter:} \emph{I can ask someone how they feel and say I am not well.}

Complete the last lesson of chapter 216: I can ask someone how they feel and say I am not well.
\end{chapteropening}

\section[\texorpdfstring{\ja{びょうき} (byōki)}{びょうき (byōki)}]{\ja{びょうき} (byōki) — an illness; a disease}

\label{lesson:JA-C216-byouki}



\begin{quote}
Before the new one: say the Japanese for a neck, then the Japanese for a cough.
\end{quote}

\subsection*{You'll want to know: \ja{びょうき}}
\textbf{\ja{びょうき}} — \emph{byōki} — \textquotedblleft{}an illness; a disease\textquotedblright{}.

\begin{quote}
\textbf{\ja{こどもは} \ja{びょうきですか。}} — \emph{kodomo wa byōki desu ka} — \textquotedblleft{}Is the child ill?\textquotedblright{}
\end{quote}

\begin{cousinweb}
\textbf{\ja{びょうき}} (\emph{byōki}) is Sino-Japanese (kango): two Chinese characters read with their on'yomi, readings borrowed from Chinese. They are \emph{byō}, \textquotedblleft{}illness\textquotedblright{} (Middle Chinese \emph{bjaengH}) and \emph{ki}, \textquotedblleft{}air, spirit\textquotedblright{} (\emph{khj+jH}) (KANJIDIC2; Middle Chinese in Baxter's transcription).
\end{cousinweb}

\begin{grammarlens}[title={What you've built}]
One more word for everyday life at A2.
\end{grammarlens}

\subsection*{Your turn}
\begin{itemize}
\raggedright
  \item \emph{Say it:} \emph{byōki}
  \item \emph{Say it:} \emph{byōki}, once more
  \item \emph{Say it:} \emph{kodomo wa byōki desu ka} — Is the child ill?
  \item \emph{Recall:} say the Japanese for a neck, then the Japanese for a cough, then say \emph{byōki} again
\end{itemize}

\begin{tcolorbox}[breakable,colback=teal!4,colframe=teal!35!black,title={Before you move on}]
What does \ja{びょうき} mean? (\textquotedblleft{}An illness; a disease\textquotedblright{}.) Say it once more.
\end{tcolorbox}
\section[\texorpdfstring{\ja{ぐあい} (guai)}{ぐあい (guai)}]{\ja{ぐあい} (guai) — condition; how one feels}

\label{lesson:JA-C216-guai}



\begin{quote}
Before the new one: say the Japanese for a cough, then the Japanese for an illness.
\end{quote}

\subsection*{You'll want to know: \ja{ぐあい}}
\textbf{\ja{ぐあい}} — \emph{guai} — \textquotedblleft{}condition; how one feels\textquotedblright{}.

\begin{quote}
\textbf{\ja{ぐあいは} \ja{いいですか。}} — \emph{guai wa ii desu ka} — \textquotedblleft{}Are you feeling well?\textquotedblright{}
\end{quote}

\begin{cousinweb}
\textbf{\ja{ぐあい}} (\emph{guai}) is built from Sino-Japanese \emph{gu}, \textquotedblleft{}tool, means\textquotedblright{} (Middle Chinese \emph{gjuH}), borrowed from Chinese, and native \emph{ai}, \textquotedblleft{}fitting\textquotedblright{}, from \emph{au}, \textquotedblleft{}to meet, fit\textquotedblright{} (JMdict, KANJIDIC2; Baxter's transcription).
\end{cousinweb}

\begin{grammarlens}[title={What you've built}]
One more word for everyday life at A2.
\end{grammarlens}

\subsection*{Your turn}
\begin{itemize}
\raggedright
  \item \emph{Say it:} \emph{guai}
  \item \emph{Say it:} \emph{guai}, once more
  \item \emph{Say it:} \emph{guai wa ii desu ka} — Are you feeling well?
  \item \emph{Recall:} say the Japanese for a cough, then the Japanese for an illness, then say \emph{guai} again
\end{itemize}

\begin{tcolorbox}[breakable,colback=teal!4,colframe=teal!35!black,title={Before you move on}]
What does \ja{ぐあい} mean? (\textquotedblleft{}Condition; how one feels\textquotedblright{}.) Say it once more.
\end{tcolorbox}
\section[\texorpdfstring{\ja{つかれる} (tsukareru)}{つかれる (tsukareru)}]{\ja{つかれる} (tsukareru) — to get tired}

\label{lesson:JA-C216-tsukareru}



\begin{quote}
Before the new one: say the Japanese for an illness, then the Japanese for condition.
\end{quote}

\subsection*{You'll want to know: \ja{つかれる}}
\textbf{\ja{つかれる}} — \emph{tsukareru} — \textquotedblleft{}to get tired\textquotedblright{}.

\begin{quote}
\textbf{\ja{まいにち} \ja{つかれますか。}} — \emph{mainichi tsukaremasu ka} — \textquotedblleft{}Do you get tired every day?\textquotedblright{}
\end{quote}

Polite \textbf{\ja{つかれます}}, \emph{tsukaremasu}; an ichidan verb.

\begin{cousinweb}
\textbf{\ja{つかれる}} (\emph{tsukareru}) is a native Japanese word (wago), inherited rather than borrowed from Chinese; its kanji was attached later as a kun'yomi, and its older origin is not settled (JMdict, KANJIDIC2).
\end{cousinweb}

\begin{grammarlens}[title={What you've built}]
One more word for everyday life at A2.
\end{grammarlens}

\subsection*{Your turn}
\begin{itemize}
\raggedright
  \item \emph{Say it:} \emph{tsukareru}
  \item \emph{Say it:} \emph{tsukareru}, once more
  \item \emph{Say it:} \emph{mainichi tsukaremasu ka} — Do you get tired every day?
  \item \emph{Recall:} say the Japanese for an illness, then the Japanese for condition, then say \emph{tsukareru} again
\end{itemize}

\begin{tcolorbox}[breakable,colback=teal!4,colframe=teal!35!black,title={Before you move on}]
What does \ja{つかれる} mean? (\textquotedblleft{}To get tired\textquotedblright{}.) Say it once more.
\end{tcolorbox}
\section[\texorpdfstring{\ja{なおる} (naoru)}{なおる (naoru)}]{\ja{なおる} (naoru) — to get better; to heal}

\label{lesson:JA-C216-naoru}



\begin{quote}
Before the new one: say the Japanese for condition, then the Japanese for to get tired.
\end{quote}

\subsection*{You'll want to know: \ja{なおる}}
\textbf{\ja{なおる}} — \emph{naoru} — \textquotedblleft{}to get better; to heal\textquotedblright{}.

\begin{quote}
\textbf{\ja{びょうきは} \ja{もう} \ja{なおりましたか。}} — \emph{byōki wa mō naorimashita ka} — \textquotedblleft{}Has your illness got better yet?\textquotedblright{}
\end{quote}

Polite \textbf{\ja{なおります}}, \emph{naorimasu}; a godan verb. Its partner \textbf{\ja{なおす}}, \emph{naosu}, means \textquotedblleft{}to fix\textquotedblright{}.

\begin{cousinweb}
\textbf{\ja{なおる}} (\emph{naoru}) is a native Japanese word (wago), inherited rather than borrowed from Chinese; its kanji was attached later as a kun'yomi, and it pairs with \emph{naosu}, \textquotedblleft{}to fix\textquotedblright{} (JMdict, KANJIDIC2).
\end{cousinweb}

\begin{grammarlens}[title={What you've built}]
One more word for everyday life at A2.
\end{grammarlens}

\subsection*{Your turn}
\begin{itemize}
\raggedright
  \item \emph{Say it:} \emph{naoru}
  \item \emph{Say it:} \emph{naoru}, once more
  \item \emph{Say it:} \emph{byōki wa mō naorimashita ka} — Has your illness got better yet?
  \item \emph{Recall:} say the Japanese for condition, then the Japanese for to get tired, then say \emph{naoru} again
\end{itemize}

\begin{tcolorbox}[breakable,colback=teal!4,colframe=teal!35!black,title={Before you move on}]
What does \ja{なおる} mean? (\textquotedblleft{}To get better; to heal\textquotedblright{}.) Say it once more.
\end{tcolorbox}
\section[\texorpdfstring{\ja{ふるえる} (furueru)}{ふるえる (furueru)}]{\ja{ふるえる} (furueru) — to shiver; to tremble}

\label{lesson:JA-C216-furueru}



\begin{quote}
Before the new one: say the Japanese for to get tired, then the Japanese for to get better.
\end{quote}

\subsection*{You'll want to know: \ja{ふるえる}}
\textbf{\ja{ふるえる}} — \emph{furueru} — \textquotedblleft{}to shiver; to tremble\textquotedblright{}.

\begin{quote}
\textbf{\ja{ねつで} \ja{からだが} \ja{ふるえますか。}} — \emph{netsu de karada ga furuemasu ka} — \textquotedblleft{}Are you shivering with fever?\textquotedblright{}
\end{quote}

Polite \textbf{\ja{ふるえます}}, \emph{furuemasu}; an ichidan verb.

\begin{cousinweb}
\textbf{\ja{ふるえる}} (\emph{furueru}) is a native Japanese word (wago), inherited rather than borrowed from Chinese; its kanji was attached later as a kun'yomi, and its older origin is not settled (JMdict, KANJIDIC2).
\end{cousinweb}

\begin{grammarlens}[title={What you've built}]
One more word for everyday life at A2.
\end{grammarlens}

\subsection*{Your turn}
\begin{itemize}
\raggedright
  \item \emph{Say it:} \emph{furueru}
  \item \emph{Say it:} \emph{furueru}, once more
  \item \emph{Say it:} \emph{netsu de karada ga furuemasu ka} — Are you shivering with fever?
  \item \emph{Recall:} say the Japanese for to get tired, then the Japanese for to get better, then say \emph{furueru} again
\end{itemize}

\begin{tcolorbox}[breakable,colback=teal!4,colframe=teal!35!black,title={Before you move on}]
What does \ja{ふるえる} mean? (\textquotedblleft{}To shiver; to tremble\textquotedblright{}.) Say it once more.
\end{tcolorbox}
